% Lösungen Einheit 4 von 4 – Ausklammern (zusammenbau v0.9)
\einheitenkopf{Einheit 4 von 4 · Ausklammern}

\erg{35}{a) die Zahl 4 \quad b) die Variable $y$ \quad c) die Zahl 2 und die Variable $a$, also $2a$ \quad d) die Zahl 3}
\erg{36}{a) In die Lücken: $2x$ und $5$, also $8x + 20 = 4 \cdot 2x + 4 \cdot 5$ \quad b) In die Lücken: $3a$ und $7$, also $6a + 14 = 2 \cdot 3a + 2 \cdot 7$ \quad c) In die Lücken: $2y$ und $3$, also $10y + 15 = 5 \cdot 2y + 5 \cdot 3$ \quad d) In die Lücken: $3b$ und $2$, also $9b + 6 = 3 \cdot 3b + 3 \cdot 2$}
\erg{37}{a) In die Klammer: $3x$ und $7$, also $9x + 21 = 3 \cdot (3x + 7)$ \quad b) In die Klammer: $2y$ und $5$, also $4y + 10 = 2 \cdot (2y + 5)$ \quad c) In die Klammer: $2a$ und $7$, also $10a + 35 = 5 \cdot (2a + 7)$ \quad d) In die Klammer: $2b$ und $3$, also $14b + 21 = 7 \cdot (2b + 3)$}
\erg{38}{a) Faktor finden: $5$; ausklammern: $5x + 10 = 5 \cdot (x + 2)$ \quad b) Faktor finden: $5$; ausklammern: $5x + 25 = 5 \cdot (x + 5)$ \quad c) Faktor finden: $5$; ausklammern: $5x + 30 = 5 \cdot (x + 6)$ \quad d) Faktor finden: $5$; ausklammern: $5x + 35 = 5 \cdot (x + 7)$ \quad e) $x \cdot (x + 5)$ \quad f) $3x \cdot (2x + 5)$ \quad g) $6 \cdot (x + 1)$ \quad h) $3 \cdot (x^2 + 2x + 5)$ \quad i) $6 \cdot (x + 5)$; Probe: $6 \cdot x + 6 \cdot 5 = 6x + 30$ \quad j) Summe der Teilflächen: $4 \cdot x + 4 \cdot 6 = 4x + 24$; Produkt mit Klammer: $4 \cdot (x + 6)$ \quad k) Ja; ausklammern: $0{,}1 \cdot 40 + 0{,}1 \cdot 25 = 0{,}1 \cdot (40 + 25)$; rechnen: $0{,}1 \cdot 65 = 6{,}50$ €. Der Rabatt beträgt 6,50 €. \quad l) Die 20 wurde nicht durch 5 geteilt, die Probe ergibt $5x + 100$. Richtig: $5x + 20 = 5 \cdot (x + 4)$.}
\erg{39}{$3a \cdot (a + 2b + 3)$; Probe: $3a \cdot a + 3a \cdot 2b + 3a \cdot 3 = 3a^2 + 6ab + 9a$}
\erg{40}{a) Das Minus ging verloren, die Probe ergibt $20x + 8$. Richtig: Faktor finden: $4$; ausklammern: $20x - 8 = 4 \cdot (5x - 2)$. \quad b) Richtig. Beide Glieder sind durch 5 teilbar, das Minus bleibt in der Klammer; Probe: $5 \cdot 2a - 5 \cdot 3 = 10a - 15$. \quad c) Nicht stimmen können die Rechnungen 2 und 3. 2: Im Term steht ein Minus, in der Klammer ein Plus. 3: Aus zwei Gliedern werden in der Klammer wieder zwei, hier fehlt die Eins.}
\erg{41}{a) Distributivgesetz: ausmultipliziert ist $6 \cdot (a + 7) = 6a + 42$, beide Terme sind gleichwertig; Probe mit $a = 2$: $6 \cdot 2 + 42 = 54$ und $6 \cdot 9 = 54$. \quad b) (1) falsch, z. B. $3x + 5$: 3 und 5 haben keinen gemeinsamen Faktor größer als 1. (2) wahr, denn Ausklammern und Ausmultiplizieren sind nach dem Distributivgesetz Umkehrungen. (3) wahr, denn jedes Glied wird durch den Faktor geteilt und bleibt stehen, notfalls als Eins. \quad c) Nein; auch die zweite 8 wird durch 8 geteilt, es bleibt die Eins: $8x - 8 = 8 \cdot (x - 1)$.}
\erg{42}{a) Zwei Rechtecke nebeneinander, beide 6 cm hoch, eines $b$ cm breit, eines 3 cm breit; zusammen $6b + 18$ cm². \quad b) Mit Klammer: $4 \cdot (x + 9)$; Klammer auflösen: $4x + 36$ \quad c) Die Summe aus einer Zahl $y$ und fünf wird verdreifacht.}
\erg{43}{a) Ja; Term mit Klammer: $15 \cdot (28 + 17)$; rechnen: $15 \cdot 45 = 675$ €, das sind 25 € weniger als 700 €. \quad b) Term mit Klammer: $3 \cdot (45 + 38)$; rechnen: $3 \cdot 83 = 249$ €. Der Einkauf kostet 249 €. \quad c) Nein; Term mit Klammer: $4 \cdot (7{,}5 + 5{,}5)$; rechnen: $4 \cdot 13 = 52$ m², das sind 2 m² mehr als 50 m².}
